\section{The Levi Civita Connection}

\subsection{Metric connections}

In this section we are concerned with connections on Riemannian manifolds.
Recall that the Euclidean connection on $\bR^n$ is defined by
$$ \bar\nabla_XY = X(Y^i)\pdv{x^i},\qquad \Gamma^k_{ij} = 0 . $$
This interacts nicely with the euclidean metric:
$$ \bar\nabla_X\iprod{Y,Z} = \iprod{\bar\nabla_XY,Z} + \iprod{Y,\bar\nabla_XZ} . $$
This can be verified easily:
$$ \bar\nabla_X\iprod{Y,Z} = \bar\nabla_X(\sum_i Y^iZ^i) = \sum_i X(Y^iZ^i) = \sum_i X(Y^i)Z^i + Y^iX(Z^i) , $$
while
$$ \iprod{\bar\nabla_XY,Z} = \iprod{X(Y^i)\pdv{x^i},Z} = \sum_i X(Y^i)Z^i $$
so the identity is easily seen to be true.

\bdefn

    Let $(M,g)$ be a Riemannian manifold, and $\nabla$ a connection on $M$.
    Then $\nabla$ is said to be {\emphcolor compatible} with $g$, or a {\emphcolor metric connection} if it satisfies
    $$ \nabla_X\iprod{Y,Z} = \iprod{\nabla_XY,Z} + \iprod{Y,\nabla_XZ} $$
    for all vector fields $X,Y,Z\in\X(M)$.

\edefn

\bprop

    Let $(M,g)$ be a Riemannian manifold, and $\nabla$ a connection on $M$.
    The following are equivalent:
    \benum
        \item $\nabla$ is compatible with $g$.
        \item $g$ is parallel with respect to $\nabla$: $\nabla g$.
        \item If $(E_i)$ is a smooth local frame, then the connection coefficients of $\nabla$ satisfy
        $$ \Gamma^\ell_{ki}g_{\ell j} + \Gamma^\ell_{kj}g_{i\ell} = E_k(g_{ij}) . $$
        \item If $V$ and $W$ are smooth vector fields along a curve $\gamma$, then
        $$ \drv t\iprod{V,W} = \iprod{D_tV,W} + \iprod{V,D_tW} . $$
        \item If $V,W$ are parallel vector fields along a curve $\gamma$, then $\iprod{V,W}$ is constant along $\gamma$.
        \item Given any smooth curve $\gamma$, every parallel transport map along $\gamma$ is a linear isometry.
        \item Given any smooth curve $\gamma$, every orthonormal basis at a point of $\gamma$ can be extended to a parallel
        orthonormal frame along $\gamma$.
    \eenum

\eprop

\bproof

    First we prove the equivalence of the first two points.
    Recall that
    $$ (\nabla g)(Y,Z,X) = (\nabla_Xg)(Y,Z) = X(g(Y,Z)) - g(\nabla_XY,Z) - g(Y,\nabla_XZ) =
    \nabla_X\iprod{Y,Z} - \iprod{\nabla_XY,Z} - \iprod{Y,\nabla_XZ} . $$
    Thus $\nabla g$ is zero iff $\nabla$ is compatible with $g$, as desired.

    Now we prove the equivalence of (2) and (3).
    Recall that $\nabla g$'s components are given by
    $$ g_{ij;k} = E_k(g_{ij}) - \Gamma^\ell_{ki}g_{ij} - \Gamma^\ell_{kj}g_{i\ell} , $$
    so $\nabla g$ is zero iff the identity holds.

    We prove now the equivalence of (1) and (4).
    If $\nabla$ is compatible with $g$, let $V,W$ be smooth vector fields along $\gamma\colon I\to M$.
    Let $t_0\in I$ and choose coordinates $(x^i)$ of a neighborhood of $\gamma(t_0)$, with $V=V^i\partial_i,W=W^j\partial_j$
    in a neighborhood of $t_0$.
    Then
    $$ \drv t\iprod{V,W} = \drv t(V^iW^j\iprod{\partial_i,\partial_j}) = (\dot V^iW^j + V^i\dot W^j)\iprod{\partial_i,\partial_j}
    + V^iW^j\drv t\iprod{\partial_i,\partial_j} . $$
    Note that for extendible $X\circ\gamma,Y\circ\gamma$,
    $$ \drv t\iprod{X\circ\gamma,Y\circ\gamma} = \drv t\parens{\iprod{X,Y}\circ\gamma} = d_{\gamma(t)}\iprod{X,Y}\circ\gamma'(t) . $$
    And since $\nabla$ is $g$-compatible:
    $$ d_{\gamma(t)}\iprod{X,Y} = \nabla\iprod{X,Y}|_{\gamma(t)} = (\iprod{\nabla X,Y} + \iprod{X,\nabla Y})|_{\gamma(t)} , $$
    since $df=\nabla f$ for all real-valued $f$, so this is equal to
    $$ \drv t\iprod{X\circ\gamma,Y\circ\gamma} = \iprod{\nabla_{\gamma'(t)}X|_{\gamma(t)},Y|_{\gamma(t)}} +
    \iprod{X|_{\gamma(t)},\nabla_{\gamma'(t)}Y|_{\gamma(t)}} . $$
    Taking $X=\partial_i$ and $Y=\nabla_j$ we have our desired result.

    Conversely if (4) holds, then it holds for extendible vector fields and so if $V,W$ are vector fields on $M$, we get the result
    by considering $V\circ\gamma,W\circ\gamma$ and recalling that
    $$ D_t(V\circ\gamma)(t) = \nabla_{\gamma'(t)}V . $$

    Now we prove (4) implies (5) implies (6) implies (7) implies (4).
    If $V,W$ are parallel along $\gamma$ then $D_tV=D_tW=0$ so $\iprod{V,W}$ has zero derivative and is thus constant.
    So (4) implies (5).

    Now assume (5), let $v_0,w_0\in T_{\gamma(t_0)}M$ and let $V,W$ be their parallel transports along $\gamma$, so $V(t_0)=v_0$
    and $W(t_0)=w_0$, $P^\gamma_{t_0t_1}v_0=V(t_1),P^\gamma_{t_0t_1}w_0=W(t_1)$.
    Because $V,W$ are parallel along $\gamma$, $\iprod{V,W}$ is constant so
    $$ \iprod{P^\gamma_{t_0t_1}v_0,P^\gamma_{t_0t_1}w_0} = \iprod{V(t_1),W(t_1)} = \iprod{V(t_0),W(t_0)} = \iprod{v_0,w_0}, $$
    so parallel transport is indeed an isometry.

    Now if $(b_i)$ is an orthonormal basis of $T_{\gamma(t)}M$, then extend $b_i$ via its parallel transport.
    Since parallel transport is an isometry, this remains an orthonormal basis.

    Finally assume (7), we will show (4).
    Let $V,W$ be smooth vector fields along $\gamma$, and let $(E_i)$ be an orthonormal frame parallel along $\gamma$ which exists by (7).
    Write $V=V^iE_i,W=W^jE_j$.
    Now $D_tV=\dot V^iE_i$ since $E_i$ is parallel, and thus
    $$ \iprod{D_tV,W} = \dot V^iW^j\iprod{E_i,E_j} = \sum_i \dot V^iW^i . $$
    And
    $$ \iprod{V,W} = \sum_i V^iW^i \implies \drv t\iprod{V,W} = \sum_i \dot V^iW^i + V^i\dot W^i = \iprod{D_tV,W} + \iprod{V,D_tW} $$
    as desired.
    \qed

\eproof

\bcoro

    Let $(M,g)$ be a Riemannian manifold, $\nabla$ a metric connection, and $\gamma\colon I\to M$ a curve.
    \benum
        \item $\norm{\gamma'(t)}$ is constant iff $D_t\gamma'(t)$ is orthogonal to $\gamma'(t)$ for all $t\in I$.
        \item If $\gamma$ is a geodesic, then $\norm{\gamma'(t)}$ is constant.
    \eenum

\ecoro

\bproof

    By the above proposition
    $$ \drv t\iprod{\gamma'(t),\gamma'(t)} = 2\iprod{\gamma'(t),D_t\gamma'(t)} , $$
    and the result follows.
    \qed

\eproof

We now turn to another nice property of the Euclidean connection.
Notice that for two vector fields $X,Y$ on Euclidean space
$$ [X,Y] = X(Y^i)\pdv{x^i} - Y(X^i)\pdv{x^i} = \bar\nabla_XY - \bar\nabla_YX . $$
This is a nice property, so we give it a name.

\bdefn

    A connection $\nabla$ on a manifold $M$ is {\emphcolor symmetric} if
    $$ \nabla_XY - \nabla_YX = [X,Y] $$
    for all vector fields $X,Y\in\X(M)$.

\edefn

Note that for a connection $\nabla$ we can define a tensor, the {\emphcolor torsion tensor}, $\tau\colon\X(M)\times\X(M)\to\X(M)$ by
$$ \tau(X,Y) = \nabla_XY - \nabla_YX - [X,Y] . $$
This is easily seen to be $C^\infty(M)$-multilinear, and thus a $(1,2)$-tensor field.
So $\nabla$ is symmetric iff its torsion is identically zero.

Notice relative to a coordinate frame $\partial_i$,
$$ \tau(\partial_i,\partial_j) = \Gamma^k_{ij}E_k - \Gamma^k_{ji}E_k $$
and thus $\tau$ vanishes iff $\Gamma^k_{ij}=\Gamma^k_{ji}$ for all $i,j,k$.
So $\nabla$ is symmetric iff $\Gamma^k_{ij}=\Gamma^k_{ji}$ relative to every coordinate frame (not necessarily every local frame).

\bthrm[title=Fundamental theorem of Riemannian geometry]

    On a Riemannian manifold $(M,g)$ there exists a unique symmetric metric connection $\nabla$.
    This connection is called the {\emphcolor Levi-Civita connection} (or {\emphcolor Riemannian connection}).

\ethrm

\bproof

    We show uniqueness first.
    If $\nabla$ is symmetric and metric, then
    $$ \eqalign{
        X\iprod{Y,Z} &= \iprod{\nabla_XY,Z} + \iprod{Y,\nabla_XZ},\cr
        Y\iprod{Z,X} &= \iprod{\nabla_YZ,X} + \iprod{Z,\nabla_YX},\cr
        Z\iprod{X,Y} &= \iprod{\nabla_ZX,Y} + \iprod{X,\nabla_ZY}.\cr
    } $$
    By symmetry we get
    $$ \eqalign{
        X\iprod{Y,Z} &= \iprod{\nabla_XY,Z} + \iprod{Y,\nabla_ZX} - \iprod{Y,[X,Z]},\cr
        Y\iprod{Z,X} &= \iprod{\nabla_YZ,X} + \iprod{Z,\nabla_XY} - \iprod{Z,[Y,X]},\cr
        Z\iprod{X,Y} &= \iprod{\nabla_ZX,Y} + \iprod{X,\nabla_YZ} - \iprod{X,[Z,Y]}.\cr
    } $$
    Adding the first two and subtracting the third yields
    $$ X\iprod{Y,Z} + Y\iprod{Z,X} - Z\iprod{X,Y} = 2\iprod{\nabla_XY,Z} + \iprod{Y,[X,Z]} + \iprod{Z,[Y,X]} - \iprod{X,[Z,Y]} . $$
    Solving for $\iprod{\nabla_XY,Z}$ yields
    $$ \iprod{\nabla_XY,Z} = \frac12\bigl(X\iprod{Y,Z} + Y\iprod{Z,X} - Z\iprod{X,Y} - \iprod{Y,[X,Z]} - \iprod{Z,[Y,X]} + \iprod{X,[Z,Y]}\bigr) . $$
    This uniquely defines $\nabla$.

    We can then use this formula to prove that the resulting connection is indeed a connection.
    By the formula above, if $(U,(x^i))$ is a coordinate chart then
    $$ \iprod{\nabla_{\partial_i}\partial_j,\partial_\ell} =
    \frac12\parens{\partial_i\iprod{\partial_j,\partial_\ell} + \partial_j\iprod{\partial_\ell,\partial_i}
    - \partial_\ell\iprod{\partial_i,\partial_j}} . $$
    Recall that $g_{ij}=\iprod{\partial_i,\partial_j}$ and $\nabla_{\partial_i}\partial_j=\Gamma^k_{ij}\partial_k$.
    So
    $$ \Gamma^k_{ij}g_{k\ell} = \frac12\parens{\partial_ig_{j\ell} + \partial_ig_{i\ell} - g_\ell g_{ij}} . $$
    Multiplying by $g$'s inverse, we see
    $$ \Gamma^k_{ij} = \frac12g^{k\ell}\parens{\partial_ig_{j\ell} + \partial_ig_{i\ell} - g_\ell g_{ij}} . $$
    This defines a connection in each chart, and by uniqueness the connections from different charts agree on their overlaps.
    Thus $\nabla$ is a well-defined connection.

    By symmetry of $g_{ij}$, we see $\Gamma^k_{ij}=\Gamma^k_{ji}$ in every coordinate domain, so $\nabla$ is symmetric.

    To show compatibility, we just show $\Gamma^\ell_{ki}g_{\ell j}+\Gamma^\ell_{kj}g_{i\ell}=\partial_k(g_{ij})$.
    This is just a direct computation.
    \qed

\eproof

The proof of the theorem gave us an explicit formula for the Levi-Civita connection, as well as explicit formulas for its
connection coefficients.
The connection coefficients of the Levi-Civita connection of $(M,g)$ are called the {\emphcolor Christoffel symbols} of $g$.

\bprop

    The Levi-Civita connection is natural: if $\phi\colon(M,g)\to(\tilde M,\tilde g)$ be an isometry of Riemannian manifolds.
    Let $\tilde\nabla$ be the Levi-Civita connection on $\tilde M$, then $\nabla=\phi^*\tilde\nabla$ is the Levi-Civita
    connection on $M$.

\eprop

\bproof

    It is sufficient to show that the pullback is symetric and metric.
    $\phi$ being an isometry means
    $$ \iprod{Y_p,Z_p} = \iprod{d\phi_pY_p,d\phi_pZ_p} = \iprod{(\phi_*Y)_{\phi(p)},(\phi_*Z)_{\phi(p)}} . $$
    Meaning $\iprod{Y,Z}=\iprod{\phi_*Y,\phi_*Z}\circ\phi$.
    Thus
    $$ \eqalign{
        X\iprod{Y,Z} &= X(\iprod{\phi_*Y,\phi_*Z}\circ\phi)\cr
            &= (\phi_*X\iprod{\phi_*Y,\phi_*Z})\circ\phi\cr
            &= (\iprod{\tilde\nabla_{\phi_*X}(\phi_*Y),\phi_*Z} + \iprod{\phi_*Y,\tilde\nabla_{\phi_*X}\phi_*Z})\circ\phi.
    } $$
    Since $\phi^{-1}$ is an isometry, by the same reason we have
    $$ \iprod{\tilde\nabla_{\phi_*X}(\phi_*Y),\phi_*Z}\circ\phi = \iprod{d\phi^{-1}_*\tilde\nabla_{\phi_*X}(\phi_*Y),Z} $$
    which is just $\iprod{\nabla_XY,Z}$ by definition, so we have the desired equality for a metric connection.

    Symmetry is due to pullbacks commuting with Lie brackets:
    $$ \eqalign{
        \nabla_XY - \nabla_YX &= \phi^{-1}_*(\tilde\nabla_{\phi_*X}(\phi_*Y) - \tilde\nabla_{\phi_*Y}(\phi_*X))\cr
            &= \phi^{-1}_*[\phi_*X,\phi_*Y]\cr
            &= [X,Y].
    } $$
    As desired.
    \qed

\eproof

\bcoro

    An isometry maps geodesics to geodesics.

\ecoro

\bproof

    If $\phi$ is an isometry, recall that as a diffeomorphism it maps $\phi^*\nabla$-geodesics to $\nabla$-geodesics.
    Since these are the Levi-Civita connections of their respective manifolds, the result follows.

\eproof

\bcoro[title=Naturality of the exponential map]

    If $\phi\colon(M,g)\to(\tilde M,\tilde g)$ is an isometry of Riemannian manifolds with Levi-Civita connections
    $\nabla,\tilde\nabla$, then for every $p\in M$ the following diagram commutes:

    \commsq{\E_p&\tilde E_{\phi p}\cr M&\tilde M}
        {d\phi_p}{\phi}{\exp_p}{\exp_{\phi(p)}}

    Where $\E_p$ is the domain of the restricted exponential $\exp_p$.

\ecoro

\bproof

    $\phi$ maps geodesics of $M$ to geodesics of $\tilde M$.
    So if $\gamma$ is a geodesic starting at $p$ with initial velocity $v\in T_pM$, then $\phi\circ\gamma$ is a geodesic
    starting at $\phi(p)$ with initial velocity $d\phi_pv$.
    Thus
    $$ \exp_{\phi p}(d\phi_pv) = \phi\circ\gamma(1) = \phi\circ\exp_p(v) $$
    as desired.
    \qed

\eproof

\bprop

    Let $(M,g)$ be Riemannian.
    Then the Levi-Civita connection $\nabla$ induces a connection on the tensor bundles compatible with the metric in the
    sense that
    $$ X\iprod{F,G} = \iprod{\nabla_XF,G} + \iprod{F,\nabla_XG} $$
    for every $X\in\X(M)$ and $F,G\in\Gamma(T^{(k,\ell)}TM)$.

\eprop

\bproof

    It is sufficient to prove this for tensor fields which are the tensors of vector and covector fields,
    $\alpha_1\otimes\cdots\alpha_{k+\ell}$.
    In such a case the result follows from the formula for the inner product of tensors and the connection on tensors.
    \qed

\eproof

\bprop

    The musical isomorphisms commute with the Levi-Civita connection.
    That is, if $F$ is a tensor field with contravariant index $i$ and $\flat$ represents lowering the index, then
    $$ \nabla(F^\flat) = (\nabla F)^\flat , $$
    and if $G$ has covariant index $i$ and $\sharp$ denotes raising the $i$th index then
    $$ \nabla(G^\sharp) = (\nabla G)^\sharp . $$

\eprop

\bproof

    Recall that $F^\flat=\tr(F\otimes g)$ where the trace is taken on the $i$th and last index of $F\otimes g$.
    Since $g$ is parallel $\nabla(F\otimes g)=\nabla(F)\otimes g+F\otimes\nabla(g)=\nabla(F)\otimes g$.
    Thus
    $$ \nabla(F^\flat) = \nabla(\tr(F\otimes g)) = \tr(\nabla(F\otimes g)) = \tr(\nabla F\otimes g) = (\nabla F)^\flat . $$
    Similar for sharp.
    \qed

\eproof

\subsection{Normal neighborhoods}

Let $(M,g)$ be a Riemannian manifold.
Recall that for $p\in M$ the restricted exponential map $\exp_p\colon\E_p\to M$ is defined on an open subset
$\E_p\subseteq T_pM$, and its differential at zero is invertible.
Thus there is an open neighborhood $V$ of the origin of $T_pM$ such that $U=\exp_p(V)$ is open and $\exp_p\colon V\to U$
is a diffeomorphism.

\bdefn

    If $U$ is a neighborhood of $p\in M$ which is the diffeomorphic image under $\exp_p$ of a {\it star-shaped\/} neighborhood
    of $0\in T_pM$ then $U$ is called a {\emphcolor normal neighborhood} of $p$.

    If $U$ is a neighborhood of $p\in M$ diffeomorphic under $\exp_p$ to an open $0\in V\subseteq T_pM$, then
    consider an open ball $B_\epsilon(0)\subseteq V$ (with respect to the Riemannian metric).
    Then $\exp_p(B_\epsilon(0))$ is a normal neighborhood of $p$, called a {\emphcolor normal ball}.
    If $\cl B_\epsilon(0)\subseteq V$, then $\exp_p(S_\epsilon(0))$, where $S_\epsilon(0)=\partial B_\epsilon(0)$,
    is called the {\emphcolor normal sphere} around $p$.

\edefn

Now if we have an orthonormal basis $(b_i)$ for $T_pM$, consider the basis isomorphism $B\colon T_pM\to\bR^n$
which maps $x^ib_i\mapsto(x^1,\dots,x^n)$.
If $U=\exp_p(V)$ is a normal neighborhood of $p$, consider the smooth coordinate map
$\phi=B^{-1}\circ(\exp_p|_V)^{-1}\colon U\to\bR^n$.
This coordinate chart $(U,\phi)$ is called the {\emphcolor Riemannian normal coordinate chart} centered at $p$,
its coordinates are the {\emphcolor Riemannian normal coordinates}.

Riemannian normal coordinates are unique in the following sense.

\bprop

    Let $(M,g)$ be a Riemannian manifold, $p\in M$, and $U$ a normal neighborhood of $p$.
    For every normal coordinate chart on $U$ centered at $p$, the coordinate basis is orthonormal at $p$;
    and for every orthonormal basis $(b_i)$ for $T_pM$ there is a unique normal coordinate chart $(x^i)$ on $U$
    such that $\partial_i|_p=b_i$ for all $i$.

    Any two normal coordinate charts $(x^i),(\tilde x^j)$ are related by
    $$ \tilde x^j = A^j_ix^i, $$
    for some constant matrix $(A^j_i)\in\orth(n)$.
    (This is not true for pseudo-Riemannian manifolds.)

\eprop

\bproof

    Let $\phi$ be a normal coordinate chart on $U$ centered at $p$ with coordinates $(x^i)$.
    This means that $\phi=B^{-1}\circ\exp_p^{-1}$ where $B\colon\bR^n\to T_pM$ is the basis isomorphism determined by
    an orthonormal basis $(b_i)$ for $T_pM$.

    Note that $d\phi^{-1}_p=d(\exp_p)_0\circ dB_0=B$ since $d(\exp_p)_0$ is the identity and $B$ is linear.
    Thus
    $$ \partial_i|_p = d\phi^{-1}_p\partial_i|_0 = B(\partial_i|_0) = b_i $$
    showing that $\partial_i|_p=b_i$ is indeed orthonormal.

    Conversely any orthonormal basis $(b_i)$ yields a normal coordinate chart $\phi$ given by $B^{-1}\circ\exp_p^{-1}$
    where $B$ is $(b_i)$'s basis isomorphism.
    The computation shows that $\partial_i|_p=b_i$.

    If $\tilde\phi=\tilde B^{-1}\circ\exp_p^{-1}$ is another such chart then
    $$ \tilde\phi\circ\phi^{-1} = \tilde B^{-1}\circ\exp_p^{-1}\circ\exp_p\circ B = \tilde B^{-1}\circ B , $$
    which is a linear isometry of $\bR^n$ (since $B$ maps an orthonormal basis to an orthonormal basis).
    Letting $A=\tilde B^{-1}\circ B$ we have the desired result.

    Notice then that two normal coordinates $(x^i)$ and $(\tilde x^i)$ are equal iff $A$ is the identity
    which is iff $\tilde\phi=\phi$.
    Thus a normal coordinate chart and normal coordinates uniquely determine one another.
    \qed

\eproof

\bprop

    Let $(M,g)$ be Riemannian, and $(U,(x^i))$ a normal coordinate chart centered at $p\in M$.
    Then,
    \benum
        \item The coordinates of $p$ are $(0,\dots,0)$,
        \item The components of the metric at $p$ are $g_{ij}=\delta_{ij}$ (or $\pm\delta_{ij}$ in the pseudo-Riemannian case),
        \item For every $v=v^i\partial_i|_p\in T_pM$, the geodesic $\gamma_v$ starting at $p$ with initial velocity $v$
        is represented in normal coordinates by the line
        $$ \gamma_v(t) = (tv^1,\dots,tv^n) , $$
        as long as $t$ is in some interval such that $\gamma_v(I)\subseteq U$.
        \item The Christoffel symbols in these coordinates vanish at $p$.
        \item All the first partial derivatives of $g_{ij}$ in these coordinates vanish at $p$.
    \eenum

\eprop

\bproof

    Point 1 is clear, point 2 is because $\partial_i|_p$ are orthonormal as per the previous proposition.
    Point 3 is because $\gamma_v(t)=\exp_p(tv)$ which in these coordinates is
    $$ \phi\circ\gamma_v(t) = \phi\circ\exp_p(tv) = B^{-1}(tv) = (tv^1,\dots,tv^n) $$
    as desired.

    Now for point 4, let $v=v^i\partial_i|_p\in T_pM$.
    Then the geodesic equation for $\gamma_v(t)=(tv^1,\dots,tv^n)$ simplifies to
    $$ \Gamma^k_{ij}(tv)v^iv^j = 0 . $$
    Evaluating at zero shows that $\Gamma^k_{ij}(0)v^iv^j=0$ for every $k$ and every $v\in T_pM$.
    In particular if we choose $v=\partial_\ell$ for some $\ell$ we see that $\Gamma^k_{\ell\ell}(0)=0$ for each $k,\ell$.
    Then choosing $v=\partial_a+\partial_b$ and $v=\partial_a-\partial_b$ we get the desired result.

    And for point 5, we have that
    $$ \Gamma^\ell_{ki}g_{\ell j} + \Gamma^\ell_{kj}g_{i\ell} = \partial_k(g_{ij}) , $$
    and the left-hand side is zero.
    \qed

\eproof

\bdefn

    Geodesics starting at $p$ and lying in a normal neighborhood of $p$ are called {\emphcolor radial geodesics},
    and the above proposition shows that in terms of normal coordinates such geodesics have the representation
    $$ \gamma_v(t) = (tv^1,\dots,tv^n) . $$

\edefn

